% Source: https://artofproblemsolving.com/community/c3h1957439_inequality
% Stable identity: aops-topic-1957439-q01
We prove the tangent-line bound
\[
\dfrac{t}{\sqrt{t+3}}\le\dfrac{7t+1}{16}\qquad(t>0),
\]
whose right side is the tangent line of $f(t)=\dfrac{t}{\sqrt{t+3}}$ at $t=1$.

Put $u=\sqrt{t+3}>\sqrt3$, so $t=u^2-3$. Multiplying by $16u>0$, the bound becomes $16(u^2-3)\le u\bigl(7(u^2-3)+1\bigr)=7u^3-20u$, i.e.
\[
7u^3-16u^2-20u+48\ge0 .
\]
This cubic factors: $(u-2)^2(7u+12)=(u^2-4u+4)(7u+12)=7u^3-16u^2-20u+48$, and $(u-2)^2(7u+12)\ge0$ for $u>0$. So the bound holds, with equality only at $u=2$, i.e. $t=1$.

Adding the bound for $t=a,b,c$:
\[
\dfrac{a}{\sqrt{a+3}}+\dfrac{b}{\sqrt{b+3}}+\dfrac{c}{\sqrt{c+3}}\le\dfrac{7(a+b+c)+3}{16}=\dfrac{21+3}{16}=\dfrac{3}{2} ,
\]
with equality exactly when $a=b=c=1$. $\blacksquare$

(Alternatively, $f''(t)=-\dfrac{t+12}{4(t+3)^2\sqrt{t+3}}<0$, so $f$ is concave and Jensen's inequality gives the same bound $3f(1)=\dfrac{3}{2}$.)
